\chapter{अंकगणित}\label{ch:g12:arith}

अंकगणित \omterm{def:g10:numbers:sets}{पूर्णांकों} का अध्ययन करता है: \omterm{def:g12:arith:divides}{विभाज्यता}, \omterm{def:g12:arith:prime}{अभाज्य} संख्याएँ, शेषफल। बहुत
समय तक इसे शुद्धतम शुद्ध गणित माना जाता रहा, और आज वही हर ऑनलाइन भुगतान की
रक्षा करता है: RSA गूढ़लेखन-पद्धति इसी अध्याय में सिद्ध होने वाली बेज़ू, गाउस
और फ़र्मा की प्रमेयों पर टिकी है।

\section{विभाज्यता और यूक्लिडीय भाग}

\begin{definition}[विभाज्यता]\label{def:g12:arith:divides}
मान लीजिए $a, b \in \Z$ है। हम कहते हैं कि $b$ $a$ को \emph{विभाजित}\index{विभाज्यता}
करता है, और $b \mid a$ लिखते हैं, यदि ऐसा $k \in \Z$ हो कि $a = kb$ हो। हम यह भी कहते हैं कि $a$
$b$ का \emph{गुणज} है।
\end{definition}

\begin{proposition}\label{prop:g12:arith:divprops}
यदि $c \mid a$ और $c \mid b$ हों, तो $c$ हर \omterm{def:g10:numbers:sets}{पूर्णांक} संयोजन $au + bv$ ($u, v \in \Z$) को \omterm{def:g12:arith:divides}{विभाजित} करता
है। यदि $a \mid b$ और $b \mid a$ हों, जहाँ $a,b \in \N$, तो $a = b$। यदि $a \mid b$ और $b \neq 0$ हों, तो $\abs a \leq \abs b$।
\end{proposition}

\begin{proof}
$a = kc$, $b = lc$ लिखिए: तब $au + bv = (ku + lv)c$। शेष बिंदु $b = ka \neq 0$ होने पर $\abs k \geq 1$ के साथ $\abs{a} = \abs{k}\,\abs{b}$ से निकलते
हैं।
\end{proof}

\begin{theorem}[यूक्लिडीय भाग]\label{thm:g12:arith:euclid}
\index{यूक्लिडीय भाग}
मान लीजिए $a \in \Z$ और $b \in \N^*$ हैं। ऐसा अद्वितीय युग्म $(q, r) \in \Z \times \N$ है कि
\[
a = bq + r \qquad\text{और}\qquad 0 \leq r < b .
\]
$q$ \emph{भागफल} है और $r$ \emph{शेषफल}।
\end{theorem}

\begin{proof}
\emph{अस्तित्व:} $a$ से अधिक न होने वाले $b$ के \omterm{def:g12:arith:divides}{गुणजों} के समुच्चय का एक
सबसे बड़ा अवयव $bq$ है (वह अरिक्त और \omterm{def:g12:seq:bounded}{ऊपर परिबद्ध} है); $r = a - bq$ रखिए। अधिकतमता से
$b(q+1) > a$, इसलिए $0 \leq r < b$।
\emph{अद्वितीयता:} यदि $0 \leq r, r' < b$ के साथ $bq + r = bq' + r'$ हो, तो $b(q - q') = r' - r$ और $\abs{r' - r} < b$: $b$ से कम
\omterm{def:g10:numbers:abs}{निरपेक्ष मान} वाला $b$ का \omterm{def:g12:arith:divides}{गुणज} $0$ ही हो सकता है, इसलिए $r = r'$ और $q = q'$।
\end{proof}

\section{सर्वांगसमताएँ}

\begin{definition}[सर्वांगसमता]\label{def:g12:arith:congruence}
मान लीजिए $n \in \N^*$ है। दो \omterm{def:g10:numbers:sets}{पूर्णांक} $a, b$ \emph{$n$ के सापेक्ष
सर्वांगसम}\index{सर्वांगसमता} हैं, और इसे $a \equiv b \pmod n$ लिखते हैं, यदि $n \mid (a - b)$ हो ---
अर्थात् यदि $a$ और $b$ को $n$ से \omterm{thm:g12:arith:euclid}{यूक्लिडीय भाग} देने पर एक ही शेषफल बचे।
\end{definition}

\begin{proposition}[संक्रियाओं के साथ संगति]\label{prop:g12:arith:congops}
यदि $a \equiv b \pmod n$ और $c \equiv d \pmod n$ हों, तो
\[
a + c \equiv b + d, \qquad
ac \equiv bd, \qquad
a^k \equiv b^k \ (k \in \N) \pmod n .
\]
\end{proposition}

\begin{proof}
$n$ $(a-b) + (c-d) = (a+c) - (b+d)$ को \omterm{def:g12:arith:divides}{विभाजित} करता है, और $ac - bd = a(c - d) + d(a - b)$ भी $n$ का \omterm{def:g12:arith:divides}{गुणज} है। घात वाला नियम
गुणनफल वाले नियम से आगमन द्वारा निकलता है।
\end{proof}

\begin{method}[$n$ के सापेक्ष घातें निकालना]\label{met:g12:arith:powers}
$a^k \bmod n$ निकालने के लिए आधार को $n$ के सापेक्ष घटाइए, फिर $\pm1$ के सर्वांगसम $a$
की कोई छोटी घात ढूँढ़िए और उससे घातांक को समेट दीजिए। जैसे $2^{100} \bmod 7$: चूँकि $2^3 = 8 \equiv 1 \pmod 7$
और $100 = 3\times33 + 1$, इसलिए
\[
2^{100} = \left(2^{3}\right)^{33} \times 2 \equiv 1^{33}\times 2 = 2 \pmod 7 .
\]
\end{method}

\section{महत्तम समापवर्तक, बेज़ू और गाउस}

\begin{definition}[महत्तम समापवर्तक]\label{def:g12:arith:gcd}
मान लीजिए $a, b$ ऐसे \omterm{def:g10:numbers:sets}{पूर्णांक} हैं जो दोनों शून्य नहीं हैं। \emph{महत्तम
समापवर्तक}\index{महत्तम समापवर्तक} $\gcd(a, b)$ वह सबसे बड़ा \omterm{def:g10:numbers:sets}{पूर्णांक} है जो $a$ और
$b$ दोनों को \omterm{def:g12:arith:divides}{विभाजित} करता है। जब $\gcd(a,b) = 1$ हो, तब $a$ और $b$ \emph{सह-अभाज्य}\index{सह-अभाज्य}
कहलाते हैं।
\end{definition}

\begin{proposition}[यूक्लिड की कलनविधि]\label{prop:g12:arith:euclidalgo}
\index{यूक्लिड की कलनविधि}
यदि $a = bq + r$ ($b \neq 0$) हो, तो $\gcd(a, b) = \gcd(b, r)$। इसलिए \omterm{thm:g12:arith:euclid}{यूक्लिडीय भाग} को बार-बार दोहराने से $\gcd(a,b)$
निकल आता है: \omterm{def:g12:arith:gcd}{महत्तम समापवर्तक} अंतिम शून्येतर शेषफल है।
\end{proposition}

\begin{proof}
$a$ और $b$ का कोई भी उभयनिष्ठ भाजक $r = a - bq$ (\cref{prop:g12:arith:divprops}) को \omterm{def:g12:arith:divides}{विभाजित} करता है, अतः वह
$b$ और $r$ का भी उभयनिष्ठ भाजक है; और इसके उलट भी, क्योंकि $a = bq + r$। दोनों
युग्मों के उभयनिष्ठ भाजक एक ही हैं, इसलिए \omterm{def:g12:arith:gcd}{महत्तम समापवर्तक} भी। कलनविधि समाप्त
हो जाती है क्योंकि शेषफल अऋणात्मक \omterm{def:g10:numbers:sets}{पूर्णांकों} का निरंतर \omterm{def:g10:functions:variations}{ह्रासमान} \omterm{def:g12:seq:sequence}{अनुक्रम} बनाते
हैं।
\end{proof}

\begin{example}\label{ex:g12:arith:euclidalgo}
$\gcd(252, 198)$: $252 = 198 + 54$; $198 = 3\times54 + 36$; $54 = 36 + 18$; $36 = 2 \times 18 + 0$। अतः $\gcd(252,198) = 18$।
\end{example}

\begin{theorem}[बेज़ू सर्वसमिका]\label{thm:g12:arith:bezout}
\index{बेज़ू सर्वसमिका}
मान लीजिए $a, b$ ऐसे \omterm{def:g10:numbers:sets}{पूर्णांक} हैं जो दोनों शून्य नहीं हैं, और $d = \gcd(a,b)$ है। ऐसे
$u, v \in \Z$ हैं कि
\[
au + bv = d .
\]
विशेष रूप से, $a$ और $b$ \omterm{def:g12:arith:gcd}{सह-अभाज्य} हैं यदि और केवल यदि किन्हीं \omterm{def:g10:numbers:sets}{पूर्णांकों}
$u, v$ के लिए $au + bv = 1$ हो।
\end{theorem}

\begin{proof}
\omterm{prop:g12:arith:euclidalgo}{यूक्लिड की कलनविधि} उल्टी चलाइए: हर शेषफल पिछले दोनों का \omterm{def:g10:numbers:sets}{पूर्णांक} संयोजन है,
और आरंभिक आँकड़े $a, b$ स्वयं अपने संयोजन हैं; नीचे से ऊपर प्रतिस्थापन करते
जाने पर अंतिम शून्येतर शेषफल $d$ $a$ और $b$ का \omterm{def:g10:numbers:sets}{पूर्णांक} संयोजन बन जाता
है। (\cref{ex:g12:arith:euclidalgo} में: $18 = 54 - 36 = 54 - (198 - 3\times54) = 4\times54 - 198 =
4(252 - 198) - 198 = 4\times252 - 5\times198$।)

तुल्यता के लिए: यदि $\gcd(a,b) = 1$ हो, तो बेज़ू $u, v$ दे देता है; इसके उलट, $a$ और
$b$ का कोई भी उभयनिष्ठ भाजक $au + bv = 1$ को \omterm{def:g12:arith:divides}{विभाजित} करता है, जिससे $\gcd(a,b) = 1$ अनिवार्य
हो जाता है।
\end{proof}

\begin{theorem}[गाउस की प्रमेयिका]\label{thm:g12:arith:gauss}
\index{गाउस की प्रमेयिका}
मान लीजिए $a, b, c \in \Z$ है। यदि $a \mid bc$ और $\gcd(a, b) = 1$ हों, तो $a \mid c$।
\end{theorem}

\begin{proof}
बेज़ू $au + bv = 1$ देता है; $c$ से गुणा कीजिए: $acu + bcv = c$। बाएँ पक्ष के दोनों पद $a$
के \omterm{def:g12:arith:divides}{गुणज} हैं (दूसरा इसलिए कि $a \mid bc$), अतः $c$ भी।
\end{proof}

\begin{corollary}\label{cor:g12:arith:coprimeprod}
यदि $a \mid c$, $b \mid c$ और $\gcd(a,b) = 1$ हों, तो $ab \mid c$।
\end{corollary}

\begin{proof}
$c = ak$ लिखिए। $b \mid ak$ और $\gcd(a,b)=1$ से गाउस $b \mid k$ देता है, मान लीजिए $k = bl$; तब $c = abl$।
\end{proof}

\section{अभाज्य संख्याएँ}

\begin{definition}[अभाज्य]\label{def:g12:arith:prime}
\omterm{def:g10:numbers:sets}{पूर्णांक} $p \geq 2$ \emph{अभाज्य}\index{अभाज्य संख्या} है यदि उसके एकमात्र धनात्मक
भाजक $1$ और $p$ हों।
\end{definition}

\begin{proposition}\label{prop:g12:arith:primedivides}
हर \omterm{def:g10:numbers:sets}{पूर्णांक} $n \geq 2$ का कोई \omterm{def:g12:arith:prime}{अभाज्य} भाजक होता है; और यदि $n$ \omterm{def:g12:arith:prime}{अभाज्य} न हो, तो
उसका कोई \omterm{def:g12:arith:prime}{अभाज्य} भाजक $\leq \sqrt n$ होता है। यदि कोई \omterm{def:g12:arith:prime}{अभाज्य} $p$ किसी गुणनफल $ab$ को
\omterm{def:g12:arith:divides}{विभाजित} करे, तो $p \mid a$ या $p \mid b$ (\emph{यूक्लिड की प्रमेयिका})।
\end{proposition}

\begin{proof}
$n$ का सबसे छोटा भाजक $d \geq 2$ \omterm{def:g12:arith:prime}{अभाज्य} है ($d$ का कोई भी उचित भाजक $n$ का
उससे छोटा भाजक होता)। यदि $n = de$ भाज्य हो और $2 \leq d \leq e$ हो, तो $d^2 \leq de = n$, इसलिए $d \leq \sqrt n$।
यूक्लिड की प्रमेयिका के लिए: यदि $p \nmid a$ हो, तो $\gcd(p, a) = 1$ ($p$ के एकमात्र भाजक
$1$ और $p$ हैं), और गाउस की प्रमेयिका $p \mid b$ दे देती है।
\end{proof}

\begin{theorem}[यूक्लिड]\label{thm:g12:arith:infiniteprimes}
\omterm{def:g12:arith:prime}{अभाज्य} संख्याएँ अनंत हैं।
\end{theorem}

\begin{proof}
\omterm{def:g12:arith:prime}{अभाज्यों} की कोई भी परिमित सूची $p_1, \dots, p_k$ दी होने पर $N = p_1 p_2 \cdots p_k + 1$ पर विचार कीजिए। कोई
\omterm{def:g12:arith:prime}{अभाज्य} $p$ $N$ को \omterm{def:g12:arith:divides}{विभाजित} करता है; पर कोई भी $p_i$ $N$ को \omterm{def:g12:arith:divides}{विभाजित} नहीं
करता (शेषफल $1$ है), इसलिए $p$ ऐसा \omterm{def:g12:arith:prime}{अभाज्य} है जो सूची में नहीं है। कोई भी
परिमित सूची \omterm{def:g12:arith:prime}{अभाज्यों} को समाप्त नहीं कर सकती।
\end{proof}

\begin{theorem}[अंकगणित की मूल प्रमेय]\label{thm:g12:arith:fta}
\index{अंकगणित की मूल प्रमेय}
हर \omterm{def:g10:numbers:sets}{पूर्णांक} $n \geq 2$ \omterm{def:g12:arith:prime}{अभाज्यों} का गुणनफल है, और यह \omterm{def:g10:algebra:expand}{गुणनखंडन} गुणनखंडों के क्रम तक
अद्वितीय है:
\[
n = p_1^{\alpha_1} p_2^{\alpha_2} \cdots p_r^{\alpha_r},
\qquad p_1 < p_2 < \dots < p_r \text{ अभाज्य},\ \alpha_i \geq 1 .
\]
\end{theorem}

\begin{proof}
\emph{अस्तित्व}, प्रबल आगमन से: \omterm{def:g12:arith:prime}{अभाज्य} $n$ स्वयं अपना \omterm{def:g10:algebra:expand}{गुणनखंडन} है; अन्यथा
$2 \leq d, e < n$ वाला $n = de$, और आगमन-परिकल्पना से दोनों गुणनखंडित हो जाते हैं।
\emph{अद्वितीयता}: मान लीजिए $p_1\cdots p_s = q_1 \cdots q_t$ (\omterm{def:g12:arith:prime}{अभाज्य}, पुनरावृत्ति की छूट के साथ)।
यूक्लिड की प्रमेयिका से $p_1$ किसी $q_j$ को \omterm{def:g12:arith:divides}{विभाजित} करता है, और \omterm{def:g12:arith:prime}{अभाज्य} होने
के कारण $p_1 = q_j$; उसे काटकर यही दोहराइए। दोनों \omterm{def:g10:algebra:expand}{गुणनखंडन} पद-दर-पद मेल खा जाते
हैं।
\end{proof}

\begin{theorem}[फ़र्मा की लघु प्रमेय]\label{thm:g12:arith:fermat}
\index{फ़र्मा की लघु प्रमेय}
मान लीजिए $p$ \omterm{def:g12:arith:prime}{अभाज्य} है और $p \nmid a$ वाला $a \in \Z$ है। तब
\[
a^{p-1} \equiv 1 \pmod p .
\]
हर $a \in \Z$ के लिए (कोई सह-अभाज्यता माने बिना) $a^p \equiv a \pmod p$।
\end{theorem}

\begin{proof}
$p$ के सापेक्ष $p - 1$ \omterm{def:g10:numbers:sets}{पूर्णांकों} $a, 2a, 3a, \dots, (p-1)a$ पर विचार कीजिए। इनमें से कोई $\equiv 0$
नहीं है (यदि $1 \leq k \leq p-1$ के साथ $p \mid ka$ हो, तो यूक्लिड की प्रमेयिका $p \mid k$ अनिवार्य
कर देती है, जो असंभव है), और वे $p$ के सापेक्ष जोड़े-जोड़े भिन्न हैं (यदि
$ka \equiv la$ हो, तो $p \mid (k - l)a$, इसलिए $p \mid k - l$, अतः $k = l$)। अतः $p$ के सापेक्ष वे किसी
क्रम में संख्याएँ $1, 2, \dots, p-1$ ही हैं। सभी सर्वांगसमताओं का गुणा करने पर:
\[
a^{p-1}\,(p-1)! \equiv (p-1)! \pmod p .
\]
चूँकि $p$ $1, \dots, p-1$ में से किसी को \omterm{def:g12:arith:divides}{विभाजित} नहीं करता, इसलिए यूक्लिड की
प्रमेयिका बार-बार लगाकर $(p-1)!$ काटा जा सकता है, जिससे $a^{p-1} \equiv 1$ बचता है। दूसरा रूप
$a$ से गुणा करने पर मिलता है (और $p \mid a$ होने पर वह तुच्छ है)।
\end{proof}

\begin{example}[गूढ़लेखन में अनुप्रयोग]\label{ex:g12:arith:rsa}
फ़र्मा की प्रमेय $n$ के सापेक्ष घातांकन को उलटने योग्य बना देती है, बशर्ते
घातांक उपयुक्त ढंग से चुने जाएँ --- और यही \emph{RSA} गूढ़लेखन-पद्धति का हृदय
है। बड़े \omterm{def:g12:arith:prime}{अभाज्य} $p, q$ और $n = pq$ के साथ $n$ तथा एक घातांक $e$ सार्वजनिक कर
दिए जाते हैं; गूढ़लेखन $x \mapsto x^e \bmod n$ है। गूढ़वाचन के लिए ऐसा घातांक $d$ चाहिए जिसके
लिए $ed \equiv 1 \pmod{(p-1)(q-1)}$ हो, और उसे केवल वही निकाल सकता है जो $p$ और $q$ जानता हो ---
और $n$ से $p, q$ वापस पाने का अर्थ है सैकड़ों अंकों वाली संख्या का \omterm{def:g10:algebra:expand}{गुणनखंडन},
जो किसी भी ज्ञात कलनविधि से उचित समय में नहीं होता।
\end{example}

\section{अभ्यास}

\begin{exercise}[$\star$]\label{exo:g12:arith:1}
$2026$ को $17$ से, और $-2026$ को $17$ से \omterm{thm:g12:arith:euclid}{यूक्लिडीय भाग} देने पर भागफल और शेषफल
निकालिए।
\end{exercise}

\begin{exercise}[$\star$]\label{exo:g12:arith:2}
$10$ के सापेक्ष $7^{100}$ का शेषफल क्या है? ($7^{100}$ का अंतिम अंक क्या है?)
\end{exercise}

\begin{exercise}[$\star$]\label{exo:g12:arith:3}
\omterm{prop:g12:arith:euclidalgo}{यूक्लिड की कलनविधि} से $\gcd(1071, 462)$ निकालिए, और ऐसे \omterm{def:g10:numbers:sets}{पूर्णांक} $u, v$ ज्ञात कीजिए कि
$1071u + 462v = \gcd(1071, 462)$ हो।
\end{exercise}

\begin{exercise}[$\star$]\label{exo:g12:arith:4}
दिखाइए कि हर $n \in \Z$ के लिए $n^2$ $4$ के सापेक्ष $0$ या $1$ के सर्वांगसम
है। इससे निकालिए कि कोई \omterm{def:g10:numbers:sets}{पूर्णांक} $\equiv 3 \pmod 4$ दो वर्गों का योग कभी नहीं होता।
\end{exercise}

\begin{exercise}[$\star\star$]\label{exo:g12:arith:5}
दिखाइए कि सभी $n \in \N$ के लिए $n(n+1)(2n+1)$ $6$ से विभाज्य है।
\end{exercise}

\begin{exercise}[$\star\star$]\label{exo:g12:arith:6}
$\Z$ में \omterm{def:g12:arith:congruence}{सर्वांगसमता} $5x \equiv 3 \pmod{11}$ हल कीजिए। (संकेत: $11$ के सापेक्ष $5$ का
प्रतिलोम ज्ञात कीजिए।)
\end{exercise}

\begin{exercise}[$\star\star$]\label{exo:g12:arith:7}
$\Z \times \Z$ में डायोफैंटीय \omterm{def:g10:algebra:equation}{समीकरण}
\[
17x - 40y = 1,
\]
हल कीजिए, फिर $17x - 40y = 6$ के सभी हल बताइए।
\end{exercise}

\begin{exercise}[$\star\star$]\label{exo:g12:arith:8}
\omterm{def:g12:arith:prime}{अभाज्य} \omterm{def:g10:algebra:expand}{गुणनखंडन} की अद्वितीयता का उपयोग करके दिखाइए कि $\sqrt2$ \omterm{ex:g10:numbers:classify}{अपरिमेय} है ($a^2 = 2b^2$
के दोनों पक्षों पर $2$ के घातांक की तुलना कीजिए)।
\end{exercise}

\begin{exercise}[$\star\star\star$]\label{exo:g12:arith:9}
मान लीजिए $p$ कोई \omterm{def:g12:arith:prime}{अभाज्य} है।
\begin{enumerate}
  \item दिखाइए कि $1 \leq k \leq p - 1$ के लिए $p$ $\dbinom{p}{k}$ को \omterm{def:g12:arith:divides}{विभाजित} करता है। (संकेत:
        $k\binom pk = p\binom{p-1}{k-1}$, \cref{exo:g12:comb:7} और गाउस की प्रमेयिका का उपयोग कीजिए।)
  \item $a \geq 0$ पर आगमन से फ़र्मा की लघु प्रमेय की एक और उपपत्ति $a^p \equiv a \pmod p$ के रूप
        में निकालिए।
\end{enumerate}
\end{exercise}

\begin{exercise}[$\star\star\star$]\label{exo:g12:arith:10}
\emph{(चीनी शेषफल समस्या.)} वे सभी \omterm{def:g10:numbers:sets}{पूर्णांक} $n$ ज्ञात कीजिए जिनके लिए
\[
n \equiv 2 \pmod 3, \qquad n \equiv 3 \pmod 5, \qquad n \equiv 2 \pmod 7 .
\]
(संकेत: पहली दो शर्तें हल कीजिए, फिर तीसरी जोड़िए; बेज़ू गुणांक सहायक होंगे।)
\end{exercise}

\section{समस्या: गुप्त कूट और जाँच-अंक}

\begin{problem}[{सप्ताहांत समस्या --- सर्वांगसमताएँ हर बारकोड और हर
क्रेडिट कार्ड की पहरेदारी करती हैं, और फ़र्मा की लघु प्रमेय संसार के रहस्यों
का ताला चलाती है}]
\label{pb:g12:arith:1}
जी.\,एच.~हार्डी ने 1940 में डींग हाँकी थी कि संख्या-सिद्धांत अनुप्रयोगों से
``अछूता'' है। अस्सी वर्ष बाद हर बारकोड की बीप, हर क्रेडिट-कार्ड भुगतान और हर
गूढ़ संदेश उन्हें झुठला देता है --- और वह भी ठीक इसी अध्याय के औज़ारों से:
सर्वांगसमताएँ (\cref{prop:g12:arith:congops}), बेज़ू प्रतिलोम (\cref{thm:g12:arith:bezout}) और फ़र्मा की लघु प्रमेय
(\cref{exo:g12:arith:9})। यह समस्या कूट जाँचती है, ताले का एक खिलौना रूप तोड़ती है, और यह
सीखती है कि असली ताला टिकता क्यों है।

\medskip
\noindent\textbf{भाग I --- \omterm{def:g12:arith:congruence}{सर्वांगसमता} में प्रवाह।}
\begin{enumerate}
  \item $2026 \bmod 7$ निकालिए; फिर $7^{100}$ का अंतिम अंक ($10$ के सापेक्ष $7$ की
        घातों का चक्र ढूँढ़िए)।
  \item तेज़ घातांकन (\cref{met:g12:arith:powers}): $5^{117} \bmod 13$ निकालिए ($5^2 \equiv -1$ से शुरू कीजिए)।
  \item $3x \equiv 5 \pmod 7$ हल कीजिए।
  \item $(97, 35)$ पर \omterm{prop:g12:arith:euclidalgo}{यूक्लिड की कलनविधि} चलाइए, पीछे प्रतिस्थापन करके ऐसे \omterm{def:g10:numbers:sets}{पूर्णांक}
        $u, v$ ज्ञात कीजिए कि $97u + 35v = 1$ हो, और $97$ के सापेक्ष $35$ का प्रतिलोम
        निकालिए।
  \item ठीक-ठीक बताइए कि $a$ $n$ के सापेक्ष कब प्रतिलोमनीय है, और
        प्रतिलोम कौन-सी प्रमेय दे देती है।
\end{enumerate}

\medskip
\noindent\textbf{भाग II --- जाँच-अंक।}
\begin{enumerate}[resume]
  \item ISBN{-}10: किसी पुस्तक-कूट के दसों अंक $d_1 \dots d_{10}$ $10d_1 + 9d_2 + \dots + 2d_9 + 1d_{10} \equiv 0
        \pmod{11}$ को संतुष्ट करने
        चाहिए। असली ISBN $0\,306\,40615\,2$ की जाँच कीजिए।
  \item सिद्ध कीजिए कि ISBN की योजना \emph{हर} एक-अंकीय त्रुटि पकड़ लेती है:
        यदि कोई एक अंक $d \not\equiv 0$ से बदल जाए, तो भारित योग $1 \leq w \leq 10$ वाले $w d$ से
        बदल जाता है --- यह कभी $\equiv 0 \pmod{11}$ क्यों नहीं हो सकता (\cref{thm:g12:arith:gauss})?
  \item सिद्ध कीजिए कि वह दो सटे हुए (भिन्न) अंकों की अदला-बदली भी पकड़ लेती
        है। फिर रचना का रहस्य समझाइए: $11$ के किस गुण ने दोनों उपपत्तियाँ
        चलाईं, और \omterm{def:g12:complex:modulus}{मापांक} $10$ के साथ क्या गड़बड़ हो सकती थी?
  \item EAN{-}13 बारकोड अंकों को $10$ के सापेक्ष $1, 3, 1, 3, \dots$ भार देते हैं। $978\,2940199\,05$
        को पूरा करने वाला जाँच-अंक निकालिए। EAN किन सटी हुई अदला-बदलियों को
        पकड़ने में \emph{चूक} जाता है? ($2(a - b) \equiv 0 \pmod{10}$ कब होता है?)
  \item क्रेडिट कार्ड लून की योजना काम में लाते हैं: दाईं ओर से हर दूसरे अंक
        को दुगुना कीजिए (दुगुना $9$ से अधिक हो तो $9$ घटा दीजिए), सब
        जोड़िए, और $10$ का \omterm{def:g12:arith:divides}{गुणज} माँगिए। जाँच-संख्या $4539\,1488\,0343\,6467$ की परीक्षा
        कीजिए।
  \item एक वाक्य में: अभाज्य मापांक ने ISBN को क्या दे दिया जो $10$ से बँधे
        EAN और लून को कभी नहीं मिल सकता?
\end{enumerate}

\medskip
\noindent\textbf{भाग III --- फ़र्मा का ताला।}
\begin{enumerate}[resume]
  \item ख़ज़ाने से पहले एक जाल: $2^{10} \bmod 341$ निकालिए, $2^{340} \bmod 341$ निष्कर्ष निकालिए ---
        और फिर $341$ का \omterm{def:g10:algebra:expand}{गुणनखंडन} कीजिए। यह उदाहरण (एक \emph{फ़र्मा
        छद्म-अभाज्य}) फ़र्मा की लघु प्रमेय को अभाज्यता की जाँच के रूप में
        काम में लाने के बारे में क्या कहता है?
  \item लघुरूप में RSA: $p = 3$, $q = 11$ लीजिए, इसलिए $n = 33$ और $(p-1)(q-1) = 20$; सार्वजनिक
        घातांक $e = 3$ है। ऐसा निजी घातांक $d$ ज्ञात कीजिए कि $3d \equiv 1 \pmod{20}$ हो
        (प्रश्न 4 की विधि)।
  \item संदेश $m = 4$ गूढ़ कीजिए: $c = m^3 \bmod 33$ निकालिए।
  \item गूढ़वाचन कीजिए: $c^d \bmod 33$ निकालिए ($c \equiv -2 \pmod{33}$ का उपयोग कीजिए) और संदेश वापस
        पाइए।
  \item गूढ़वाचन सदा काम क्यों करता है: दिखाइए कि $3$ के सापेक्ष और $11$
        के सापेक्ष, दोनों जगह $m^{21} \equiv m$ (हर संसार में फ़र्मा की लघु प्रमेय), और
        $33$ के सापेक्ष निष्कर्ष निकालिए (\cref{thm:g12:arith:gauss} दोनों सर्वांगसमताओं को जोड़
        देता है)। $21 = ed$ का विशेष रूप $1 + 20k$ कहाँ काम आया?
  \item ताले की सुरक्षा: $n$ और $e$ सब जानते हैं; $d$ वापस पाने के लिए
        $(p-1)(q-1)$ चाहिए, अर्थात् $n$ के गुणनखंड। हमारा $33$ तो देखते ही
        गुणनखंडित हो जाता है --- फिर वही योजना छह सौ अंकों वाले $n$ के साथ
        संसार के बैंकों की रक्षा कैसे करती है? (गुणा करने और \omterm{def:g10:algebra:expand}{गुणनखंडन} करने
        की विषमता पर एक वाक्य।)
\end{enumerate}

\medskip
\noindent\textbf{भाग IV --- चिरपरिचित।}
\begin{enumerate}[resume]
  \item सैनिकों की पुरानी चीनी गिनती (\cref{exo:g12:arith:10} से तुलना कीजिए): सैनिकों की
        संख्या $3$ की पंक्तियों में लगाने पर शेषफल $2$ और $5$ की
        पंक्तियों में लगाने पर शेषफल $3$ छोड़ती है। सभी संभव संख्याएँ
        ज्ञात कीजिए, और समझाइए कि उत्तर $15$ के सापेक्ष अद्वितीय क्यों है।
  \item अंततः एक-पंक्ति की उपपत्तियाँ: $10 \equiv 1 \pmod 9$ से सिद्ध कीजिए कि हर संख्या
        $9$ के सापेक्ष अपने अंकों के योग के सर्वांगसम होती है; और $10 \equiv -1 \pmod{11}$
        से $11$ के लिए एकांतर-योग वाला नियम निकालिए। (माध्यमिक विद्यालय
        खंड ने इन्हें स्पष्ट बीजगणित से सिद्ध किया था --- अब इस संक्षेप की
        सराहना कीजिए।)
  \item समापन --- बारकोड के सामने हार्डी: अध्याय का औज़ार-बक्सा दोहराइए
        (\omterm{def:g12:arith:congruence}{सर्वांगसमता} का अंकगणित, बेज़ू प्रतिलोम, फ़र्मा की लघु प्रमेय,
        सह-अभाज्य मापांकों को जोड़ना) और बताइए कि इस समस्या में हर औज़ार
        कहाँ फ़िट बैठा; फिर ``अछूता'' पर आधुनिक फ़ैसला सुनाइए।
\end{enumerate}
\end{problem}
